% Lecture notes with algomotion figures. Compile with shell escape:
%   latexmk -pdf -shell-escape notes.tex     (or -pdflua, -pdfxe)
% The figures are kept in _algomotion/, so later runs need no shell escape.
\documentclass[11pt]{article}
\usepackage{algomotion}
\usepackage{amsmath}

\begin{document}

\section*{Shortest paths}

Dijkstra's algorithm settles one node at a time. Once $B$ and $C$ are
settled in the textbook example, the edge from $C$ lowers the distance
to $F$ from $14$ to $11$. The graph is written as a list of weighted
edges, and the algorithm starts at its first node, $A$:

\begin{center}
  \algomotion[step=improve#6]{dijkstra}{A-B 7, A-C 9, B-C 10, B-D 15, C-D 11,
    D-E 6, A-F 14, C-F 2, E-F 9}
\end{center}

The input can also come from a file, here drawn in the print theme at a
given width:

\begin{center}
  \algomotionfile[algorithm=dijkstra, step=-1, theme=print, width=0.55\linewidth]{graph.json}
\end{center}

\section*{Sorting}

Bubble sort swaps neighbours that are out of order; this is its second
swap. Figures of one run share one scale.

\begin{center}
  \algomotion[step=swap#2]{bubble-sort}{5 1 4 2 8}
  \par\medskip
  \algomotion[step=-1]{bubble-sort}{5 1 4 2 8}
\end{center}

\section*{Deleting from a binary search tree}

A document made beforehand, for example with
\texttt{algomotion run bst.delete} or by a model, is drawn with
\verb|\algomotionfile|. Deleting $3$ copies its successor $4$ into its node:

\begin{center}
  \algomotionfile[step=delete-3/copy]{bst-delete.json}
\end{center}

\end{document}
